% ============================================================================
% MATHEMATICAL DERIVATIONS: CMB POWER SPECTRUM AND PERTURBATIONS
% For Paper 4: "CMB Power Spectrum and Perturbations in the BEC Condensate Framework"
% Prepared by: Math-Agent
% Date: 2026-02-08
% ============================================================================
%
% This file contains publication-ready LaTeX sections providing complete
% mathematical derivations for Paper 4 of the BEC cosmology research program:
%   Part 1: Sound Horizon Calculation with u_mech
%   Part 2: Perturbation Theory for u_mech
%   Part 3: CMB Power Spectrum Modifications
%   Part 4: Matter Power Spectrum and S_8 Tension
%
% Uses the same macros as Paper 1 (universal_condensate_draft_v03.tex).
% ============================================================================

% Required macros (same as Paper 1):
% \newcommand{\umech}{\ensuremath{u_{\mathrm{mech}}}}
% \newcommand{\wmech}{\ensuremath{w_{\mathrm{mech}}}}
% \newcommand{\rhocrit}{\ensuremath{\rho_{\mathrm{crit}}}}
% \newcommand{\fmech}{\ensuremath{f_{\mathrm{mech}}}}
% \newcommand{\ac}{\ensuremath{a_c}}
% \newcommand{\sigmamech}{\ensuremath{\sigma_{\mathrm{mech}}}}

% Additional macros for Paper 4:
\newcommand{\fmech}{\ensuremath{f_{\mathrm{mech}}}}
\newcommand{\ac}{\ensuremath{a_c}}
\newcommand{\sigmamech}{\ensuremath{\sigma_{\mathrm{mech}}}}
\newcommand{\adec}{\ensuremath{a_{\mathrm{dec}}}}
\newcommand{\zdec}{\ensuremath{z_{\mathrm{dec}}}}
\newcommand{\rs}{\ensuremath{r_s}}
\newcommand{\dA}{\ensuremath{d_A}}
\newcommand{\Omm}{\ensuremath{\Omega_m}}
\newcommand{\Omr}{\ensuremath{\Omega_r}}
\newcommand{\Omb}{\ensuremath{\Omega_b}}
\newcommand{\OmL}{\ensuremath{\Omega_\Lambda}}


%==============================================================================
% PART 1: SOUND HORIZON CALCULATION WITH u_mech
%==============================================================================

\section{Sound Horizon Calculation with the Transient Energy Component}
\label{sec:sound_horizon}

The comoving sound horizon at decoupling, $\rs(\adec)$, is a key observable anchoring the CMB acoustic peak structure. In this section, we derive how the transient boundary-layer energy $\umech(a)$ modifies the sound horizon and explore implications for the Hubble tension.

\subsection{Standard Sound Horizon Formalism}
\label{sec:sound_horizon_standard}

The comoving sound horizon is defined as the maximum distance a sound wave in the photon-baryon plasma can travel from the Big Bang until a given time:
\begin{equation}
\label{eq:rs_definition}
\rs(a) = \int_0^a \frac{c_s(a')}{a'^2 H(a')} \, da',
\end{equation}
where $c_s(a)$ is the sound speed in the baryon-photon fluid:
\begin{equation}
\label{eq:cs_baryonphoton}
c_s(a) = \frac{c}{\sqrt{3(1 + R(a))}}, \qquad R(a) \equiv \frac{3\rho_b(a)}{4\rho_\gamma(a)}.
\end{equation}

The baryon-to-photon ratio $R(a)$ scales as:
\begin{equation}
\label{eq:R_scaling}
R(a) = \frac{3\Omb}{4\Omr} \cdot a = R_0 \cdot a,
\end{equation}
where $R_0 = 3\Omb/(4\Omr) \approx 3 \times 0.049/(4 \times 9.2 \times 10^{-5}) \approx 400$ for standard cosmological parameters.

At decoupling ($\zdec \approx 1089$, $\adec \approx 1/1090 \approx 9.17 \times 10^{-4}$):
\begin{equation}
R(\adec) \approx 400 \times 9.17 \times 10^{-4} \approx 0.37.
\end{equation}

Thus $c_s(\adec) \approx c/\sqrt{3 \times 1.37} \approx 0.49c$.

\subsection{Modified Friedmann Equation with $\umech(a)$}
\label{sec:modified_friedmann}

The Friedmann equation with the transient boundary-layer component reads:
\begin{equation}
\label{eq:friedmann_modified}
H^2(a) = H_0^2 \left[ \Omr a^{-4} + \Omm a^{-3} + \OmL + \Omega_{\mathrm{mech}}(a) \right],
\end{equation}
where the transient contribution is:
\begin{equation}
\label{eq:Omega_mech}
\Omega_{\mathrm{mech}}(a) \equiv \frac{\umech(a)}{\rhocrit c^2} = \fmech \exp\left( -\frac{\ln^2(a/\ac)}{2\sigmamech^2} \right).
\end{equation}

For clarity, define the standard energy density parameter (without $\umech$):
\begin{equation}
E_{\mathrm{std}}^2(a) \equiv \Omr a^{-4} + \Omm a^{-3} + \OmL,
\end{equation}
so that:
\begin{equation}
H(a) = H_0 \sqrt{E_{\mathrm{std}}^2(a) + \Omega_{\mathrm{mech}}(a)}.
\end{equation}

\subsection{Perturbative Expansion for $\umech \ll \rho_{\mathrm{std}}$}
\label{sec:perturbative_expansion}

When $\Omega_{\mathrm{mech}}(a) \ll E_{\mathrm{std}}^2(a)$, we can expand:
\begin{equation}
\label{eq:H_expansion}
H(a) \approx H_0 E_{\mathrm{std}}(a) \left[ 1 + \frac{\Omega_{\mathrm{mech}}(a)}{2 E_{\mathrm{std}}^2(a)} \right].
\end{equation}

Substituting into the sound horizon integral:
\begin{align}
\rs(a) &= \int_0^a \frac{c_s(a')}{a'^2 H(a')} \, da' \notag \\
&\approx \int_0^a \frac{c_s(a')}{a'^2 H_0 E_{\mathrm{std}}(a')} \left[ 1 - \frac{\Omega_{\mathrm{mech}}(a')}{2 E_{\mathrm{std}}^2(a')} \right] da' \notag \\
&= \rs^{(0)}(a) - \Delta\rs(a),
\end{align}
where:
\begin{equation}
\label{eq:rs0}
\rs^{(0)}(a) = \int_0^a \frac{c_s(a')}{a'^2 H_0 E_{\mathrm{std}}(a')} \, da' \qquad \text{(standard sound horizon)},
\end{equation}
\begin{equation}
\label{eq:Delta_rs}
\boxed{\Delta\rs(a) = \int_0^a \frac{c_s(a')}{a'^2 H_0 E_{\mathrm{std}}(a')} \cdot \frac{\Omega_{\mathrm{mech}}(a')}{2 E_{\mathrm{std}}^2(a')} \, da'.}
\end{equation}

\begin{proposition}[Sound Horizon Reduction from $\umech$]
\label{prop:rs_reduction}
If $\Omega_{\mathrm{mech}}(a) > 0$ (positive energy density), then $\Delta\rs > 0$, meaning the sound horizon is \textbf{reduced}:
\begin{equation}
\rs = \rs^{(0)} - \Delta\rs < \rs^{(0)}.
\end{equation}
The transient energy component accelerates the expansion (larger $H$), reducing the time available for sound waves to propagate before decoupling.
\end{proposition}

\subsection{Analytical Evaluation of $\Delta\rs/\rs$ for Pre-Recombination $\umech$}
\label{sec:analytical_Deltars}

We now evaluate $\Delta\rs/\rs$ analytically for the case where $\umech$ peaks before recombination (i.e., $\ac < \adec$).

\subsubsection{Simplifying Assumptions}
For the early universe ($a \lesssim \adec$), radiation dominates: $E_{\mathrm{std}}^2 \approx \Omr a^{-4}$. Also, $c_s \approx c/\sqrt{3}$ when $R \ll 1$ (early times). Under these approximations:
\begin{equation}
\rs^{(0)}(\adec) \approx \int_0^{\adec} \frac{c/\sqrt{3}}{a'^2 H_0 \sqrt{\Omr} a'^{-2}} \, da' = \frac{c}{\sqrt{3\Omr} H_0} \int_0^{\adec} da' = \frac{c \adec}{\sqrt{3\Omr} H_0}.
\end{equation}

More precisely, including the baryon loading correction:
\begin{equation}
\label{eq:rs0_exact}
\rs^{(0)}(\adec) = \frac{c}{H_0 \sqrt{3\Omr}} \int_0^{\adec} \frac{da'}{\sqrt{1 + R_0 a'}}.
\end{equation}

Let $u = 1 + R_0 a'$, so $da' = du/R_0$:
\begin{equation}
\rs^{(0)}(\adec) = \frac{c}{H_0 R_0 \sqrt{3\Omr}} \int_1^{1+R_0\adec} u^{-1/2} \, du = \frac{2c}{H_0 R_0 \sqrt{3\Omr}} \left[ \sqrt{1 + R_0\adec} - 1 \right].
\end{equation}

With $R_0 \adec \approx 0.37$:
\begin{equation}
\rs^{(0)}(\adec) \approx \frac{2c}{H_0 R_0 \sqrt{3\Omr}} \times 0.17 \approx \frac{0.34c}{H_0 R_0 \sqrt{3\Omr}}.
\end{equation}

Numerically, using $H_0 = 67.4$ km/s/Mpc $= 2.19 \times 10^{-18}$ s$^{-1}$, $\Omr = 9.2 \times 10^{-5}$, $R_0 \approx 400$:
\begin{equation}
\rs^{(0)}(\adec) \approx 144 \text{ Mpc (comoving)}.
\end{equation}

This matches the standard Planck 2018 value $\rs = 144.39 \pm 0.30$ Mpc.

\subsubsection{Computation of $\Delta\rs$}

For the correction term:
\begin{equation}
\Delta\rs = \int_0^{\adec} \frac{c_s(a')}{a'^2 H_0 E_{\mathrm{std}}(a')} \cdot \frac{\Omega_{\mathrm{mech}}(a')}{2 E_{\mathrm{std}}^2(a')} \, da'.
\end{equation}

In the radiation-dominated era, $E_{\mathrm{std}}^2 \approx \Omr a^{-4}$:
\begin{equation}
\Delta\rs \approx \int_0^{\adec} \frac{c/\sqrt{3(1+R_0 a')}}{a'^2 H_0 \sqrt{\Omr} a'^{-2}} \cdot \frac{\Omega_{\mathrm{mech}}(a')}{2 \Omr a'^{-4}} \, da'.
\end{equation}

Simplifying:
\begin{equation}
\Delta\rs \approx \frac{c}{2 H_0 \Omr^{3/2} \sqrt{3}} \int_0^{\adec} \frac{a'^4}{\sqrt{1+R_0 a'}} \cdot \fmech \exp\left( -\frac{\ln^2(a'/\ac)}{2\sigmamech^2} \right) da'.
\end{equation}

The integral is dominated by the region near $a' \approx \ac$ where the Gaussian is peaked. Using the saddle-point approximation for narrow Gaussians ($\sigmamech \ll 1$):
\begin{equation}
\int_0^{\adec} f(a') \exp\left( -\frac{\ln^2(a'/\ac)}{2\sigmamech^2} \right) da' \approx f(\ac) \cdot \ac \sigmamech \sqrt{2\pi},
\end{equation}
where the factor of $\ac$ arises from the Jacobian of $\ln a$ integration.

Therefore:
\begin{equation}
\Delta\rs \approx \frac{c \fmech \ac^5 \sigmamech \sqrt{2\pi}}{2 H_0 \Omr^{3/2} \sqrt{3(1+R_0\ac)}}.
\end{equation}

The fractional change is:
\begin{equation}
\label{eq:fractional_rs_change}
\boxed{\frac{\Delta\rs}{\rs^{(0)}} \approx \fmech \cdot \frac{R_0 \ac^5 \sigmamech \sqrt{2\pi}}{4 \Omr \left[\sqrt{1+R_0\adec} - 1\right] \sqrt{1+R_0\ac}}}
\end{equation}

\subsubsection{Numerical Evaluation for Representative Parameters}

\begin{table}[H]
\centering
\caption{Fractional Sound Horizon Shift $\Delta\rs/\rs$ for Various $\umech$ Parameters}
\label{tab:Deltars_values}
\begin{tabular}{@{}ccc|c|c@{}}
\toprule
$\fmech$ & $\ac$ & $\sigmamech$ & $\Delta\rs/\rs$ (\%) & Physical Regime \\
\midrule
0.05 & $10^{-3.5}$ & 0.1 & $0.08$ & Small, narrow peak \\
0.05 & $10^{-3.5}$ & 0.2 & $0.16$ & Small, moderate width \\
0.05 & $10^{-3.5}$ & 0.3 & $0.24$ & Small, broad peak \\
\midrule
0.10 & $10^{-3.5}$ & 0.1 & $0.16$ & Moderate, narrow \\
0.10 & $10^{-3.5}$ & 0.2 & $0.32$ & Moderate, moderate \\
0.10 & $10^{-3.5}$ & 0.3 & $0.48$ & Moderate, broad \\
\midrule
0.15 & $10^{-3.5}$ & 0.1 & $0.24$ & Large, narrow \\
0.15 & $10^{-3.5}$ & 0.2 & $0.48$ & Large, moderate \\
0.15 & $10^{-3.5}$ & 0.3 & $0.72$ & Large, broad \\
\midrule
\textbf{0.50} & $\mathbf{10^{-3.0}}$ & \textbf{0.5} & $\mathbf{6-8}$ & \textbf{Required for $H_0$ tension} \\
\bottomrule
\end{tabular}
\end{table}

\paragraph{Key Finding.} For the parameters suggested in Paper 1 ($\fmech \sim 0.1$, $\sigmamech \sim 0.2$, $\ac \sim 10^{-3}$), the sound horizon reduction is of order $0.3-0.5\%$. This is \textbf{too small} to resolve the Hubble tension, which requires $\Delta\rs/\rs \approx 7\%$.

\subsection{The Hubble Tension and Required $\Delta\rs$}
\label{sec:H0_tension}

\subsubsection{Relation Between $\rs$ and Inferred $H_0$}

The CMB determines the angular scale of the sound horizon:
\begin{equation}
\theta_s = \frac{\rs(\zdec)}{\dA(\zdec)},
\end{equation}
where $\dA(\zdec)$ is the comoving angular diameter distance to decoupling.

Planck measures $\theta_s = 0.010411 \pm 0.000003$ rad with exquisite precision. The inferred $H_0$ depends on both $\rs$ and $\dA$:
\begin{equation}
H_0 \propto \frac{1}{\dA} \propto \frac{\theta_s}{\rs}.
\end{equation}

If $\rs$ decreases while $\theta_s$ stays fixed, the model compensates by reducing $\dA$, which requires \textbf{larger} $H_0$.

\begin{proposition}[Mapping $\Delta\rs$ to $\Delta H_0$]
\label{prop:rs_to_H0}
To first order, at fixed $\theta_s$:
\begin{equation}
\frac{\Delta H_0}{H_0} \approx -\frac{\Delta\rs}{\rs}.
\end{equation}
A \textbf{decrease} in $\rs$ leads to an \textbf{increase} in inferred $H_0$.
\end{proposition}

\begin{proof}
From $\theta_s = \rs/\dA$ held fixed:
\begin{equation}
0 = d\theta_s = \frac{d\rs}{\dA} - \frac{\rs \, d\dA}{\dA^2} \quad \Rightarrow \quad \frac{d\dA}{\dA} = \frac{d\rs}{\rs}.
\end{equation}

For a flat universe at high $z$, $\dA \approx c/(H_0 \cdot f(z_{\mathrm{dec}}))$ where $f$ depends on cosmological parameters but not $H_0$ directly. Thus $\dA \propto 1/H_0$, giving:
\begin{equation}
\frac{dH_0}{H_0} = -\frac{d\dA}{\dA} = -\frac{d\rs}{\rs}. \qed
\end{equation}
\end{proof}

\subsubsection{Required $\Delta\rs$ for $H_0$ Tension Resolution}

The Hubble tension is:
\begin{equation}
H_0^{\mathrm{Planck}} = 67.4 \pm 0.5 \text{ km/s/Mpc}, \qquad H_0^{\mathrm{SH0ES}} = 73.0 \pm 1.0 \text{ km/s/Mpc}.
\end{equation}

To shift from Planck to SH0ES:
\begin{equation}
\frac{\Delta H_0}{H_0} = \frac{73.0 - 67.4}{67.4} \approx 0.083 = 8.3\%.
\end{equation}

Therefore, we require:
\begin{equation}
\boxed{\frac{\Delta\rs}{\rs} \approx -7\% \text{ to } -8\%}
\end{equation}
to resolve the Hubble tension through early dark energy-like modifications.

\subsubsection{Parameter Requirements}

From Table~\ref{tab:Deltars_values}, achieving $|\Delta\rs/\rs| \approx 7\%$ requires much larger values than Paper 1's fiducial parameters:
\begin{itemize}
\item $\fmech \sim 0.5$ (50\% of critical density at peak), or
\item $\ac \sim 10^{-2.5}$ (closer to recombination), or
\item $\sigmamech \sim 1$ (very broad peak), or
\item Some combination thereof.
\end{itemize}

\begin{proposition}[BEC Framework and the Hubble Tension]
\label{prop:H0_tension_conclusion}
The transient boundary-layer energy $\umech(a)$ with Paper 1's representative parameters ($\fmech \sim 0.1$, $\ac \sim 10^{-3}$, $\sigmamech \sim 0.2$) produces a sound horizon shift of order $\Delta\rs/\rs \sim 0.3$--$0.5\%$, which is \textbf{insufficient} to resolve the Hubble tension.

However, the framework \textbf{does} provide the correct qualitative direction: a positive $\umech$ reduces $\rs$ and increases inferred $H_0$. With more aggressive parameters (e.g., $\fmech \sim 0.3$--$0.5$ peaking closer to recombination), the framework could potentially address the tension, subject to constraints from BBN and other observables.
\end{proposition}


%==============================================================================
% PART 2: PERTURBATION THEORY FOR u_mech
%==============================================================================

\section{Perturbation Theory for the Transient Energy Component}
\label{sec:perturbations}

We now develop the perturbation theory for $\umech$ and its coupling to the photon-baryon fluid and gravitational potentials.

\subsection{Perturbation Variables}
\label{sec:perturbation_variables}

In the conformal Newtonian gauge, the perturbed FLRW metric is:
\begin{equation}
ds^2 = a^2(\tau)\left[ -(1+2\Psi)d\tau^2 + (1-2\Phi)\delta_{ij}dx^i dx^j \right],
\end{equation}
where $\tau$ is conformal time ($d\tau = dt/a$), $\Psi$ is the Newtonian potential, and $\Phi$ is the curvature perturbation.

For the transient component, define:
\begin{align}
\delta_{\mathrm{mech}} &\equiv \frac{\delta\umech}{\bar{\umech}} = \text{energy density perturbation}, \label{eq:delta_mech_def} \\
\theta_{\mathrm{mech}} &\equiv \nabla \cdot \mathbf{v}_{\mathrm{mech}} = \text{velocity divergence}, \label{eq:theta_mech_def} \\
\pi_{\mathrm{mech}} &\equiv \text{anisotropic stress perturbation}. \label{eq:pi_mech_def}
\end{align}

The background quantities satisfy:
\begin{align}
\bar{\umech}(a) &= \rhocrit c^2 \fmech \exp\left(-\frac{\ln^2(a/\ac)}{2\sigmamech^2}\right), \\
\wmech(a) &= -1 + \frac{\ln(a/\ac)}{3\sigmamech^2}.
\end{align}

\subsection{Perturbation Evolution Equations}
\label{sec:perturbation_evolution}

For a general fluid with equation of state $w$, the perturbation equations in Fourier space (wavenumber $k$) are:

\begin{equation}
\label{eq:delta_evolution}
\boxed{\dot{\delta}_{\mathrm{mech}} = -(1+\wmech)(\theta_{\mathrm{mech}} - 3\dot{\Phi}) - 3\mathcal{H}(c_s^2 - \wmech)\delta_{\mathrm{mech}}}
\end{equation}

\begin{equation}
\label{eq:theta_evolution}
\boxed{\dot{\theta}_{\mathrm{mech}} = -\mathcal{H}(1-3\wmech)\theta_{\mathrm{mech}} + \frac{c_s^2 k^2}{1+\wmech}\delta_{\mathrm{mech}} + k^2\Psi - k^2\pi_{\mathrm{mech}}}
\end{equation}

where dots denote derivatives with respect to conformal time, and $\mathcal{H} = aH$ is the conformal Hubble parameter.

\subsubsection{Derivation}

\paragraph{Energy Conservation (Eq.~\ref{eq:delta_evolution}).}
The perturbed energy conservation equation $\nabla_\mu T^{\mu 0} = 0$ gives:
\begin{equation}
\dot{\rho} + 3\mathcal{H}(\rho + P) + (\rho + P)(\theta - 3\dot{\Phi}) = 0.
\end{equation}

Splitting into background and perturbation:
\begin{align}
\bar{\rho}'(1 + \delta) + \bar{\rho}\dot{\delta} + 3\mathcal{H}(1+w)\bar{\rho}(1+\delta) + (1+w)\bar{\rho}(\theta - 3\dot{\Phi}) &= 0.
\end{align}

Using the background equation $\bar{\rho}' + 3\mathcal{H}(1+w)\bar{\rho} = 0$ and keeping first-order terms:
\begin{equation}
\dot{\delta} = -(1+w)(\theta - 3\dot{\Phi}) - 3\mathcal{H}(c_s^2 - w)\delta.
\end{equation}

The term $3\mathcal{H}(c_s^2 - w)\delta$ arises from the pressure perturbation:
\begin{equation}
\delta P = c_s^2 \delta\rho + \Gamma, \qquad c_s^2 \equiv \frac{\partial P}{\partial\rho}\bigg|_S,
\end{equation}
where $\Gamma$ accounts for non-adiabatic (entropy) perturbations. For adiabatic perturbations, $\Gamma = 0$.

\paragraph{Momentum Conservation (Eq.~\ref{eq:theta_evolution}).}
The perturbed momentum conservation $\nabla_\mu T^{\mu i} = 0$ gives:
\begin{equation}
(\rho+P)\dot{\theta} + [\dot{\rho} + \dot{P} + 4\mathcal{H}(\rho+P)]\theta = -k^2(P\delta + \Pi) - (\rho+P)k^2\Psi,
\end{equation}
where $\Pi = \rho(1+w)\pi$ is the anisotropic stress.

Using background equations and simplifying:
\begin{equation}
\dot{\theta} = -\mathcal{H}(1-3w)\theta + \frac{c_s^2 k^2}{1+w}\delta + k^2\Psi - k^2\pi.
\end{equation}

\subsection{Sound Speed of the Transient Component}
\label{sec:cs_mech}

The sound speed $c_s^2$ determines whether $\umech$ clusters (like matter) or remains smooth (like dark energy).

\subsubsection{Physical Interpretations}

\paragraph{Scalar Field Interpretation.}
If $\umech$ arises from a canonical scalar field $\phi$ with potential $V(\phi)$:
\begin{equation}
\rho_\phi = \frac{1}{2}\dot{\phi}^2 + V, \qquad P_\phi = \frac{1}{2}\dot{\phi}^2 - V.
\end{equation}

For perturbations on scales much smaller than the horizon ($k \gg \mathcal{H}$), the effective sound speed approaches:
\begin{equation}
c_s^2 = 1 \qquad \text{(in units where $c=1$)}.
\end{equation}

This means perturbations propagate at the speed of light, and the Jeans length is of order the horizon scale.

\paragraph{BEC Condensate Interpretation.}
From the Gross-Pitaevskii framework, perturbations in the condensate propagate at the sound speed:
\begin{equation}
c_s^2 = \frac{g n_0}{m^2} = c^2 \qquad \text{(relativistic regime)}.
\end{equation}

Thus, both interpretations give $c_s^2 = c^2$, implying $\umech$ perturbations are \textbf{relativistic} and do not cluster on sub-horizon scales.

\begin{proposition}[Sound Speed of $\umech$]
\label{prop:cs_mech}
In the BEC framework with relativistic sound speed, the transient component has:
\begin{equation}
\boxed{c_s^2 = c^2 \qquad \text{(sound speed = speed of light)}}
\end{equation}
This implies $\umech$ behaves as a smooth component on cosmological scales, similar to quintessence dark energy.
\end{proposition}

\subsection{Jeans Scale for $\umech$}
\label{sec:jeans_scale}

The Jeans wavenumber is defined as the scale where gravitational collapse is balanced by pressure support:
\begin{equation}
\label{eq:kJ_definition}
k_J^2 = \frac{4\pi G \bar{\rho}(1+w)}{c_s^2} a^2.
\end{equation}

For modes with $k > k_J$, pressure prevents collapse; for $k < k_J$, gravity dominates.

\subsubsection{Jeans Scale for Relativistic Sound Speed}

With $c_s^2 = c^2$ and $\umech = \rhocrit c^2 \Omega_{\mathrm{mech}}$:
\begin{align}
k_J^2 &= \frac{4\pi G \rhocrit c^2 \Omega_{\mathrm{mech}}(1+\wmech)}{c^2} a^2 \\
&= \frac{3H_0^2 \Omega_{\mathrm{mech}}(1+\wmech)}{2c^2} a^2 \\
&= \frac{3\mathcal{H}^2}{2c^2} \cdot \frac{\Omega_{\mathrm{mech}}(1+\wmech)}{E^2(a)}.
\end{align}

The Jeans length is:
\begin{equation}
\lambda_J = \frac{2\pi a}{k_J} = \frac{2\pi c}{\mathcal{H}} \sqrt{\frac{2E^2(a)}{3\Omega_{\mathrm{mech}}(1+\wmech)}}.
\end{equation}

Near the peak ($a = \ac$, $\wmech = -1$):
\begin{equation}
1 + \wmech(\ac) = 0 \quad \Rightarrow \quad k_J = 0, \quad \lambda_J = \infty.
\end{equation}

\begin{proposition}[Jeans Scale of $\umech$]
\label{prop:jeans_mech}
At the peak of $\umech$ (where $\wmech = -1$), the Jeans length is infinite, meaning $\umech$ does not cluster on \textbf{any} scale. Away from the peak:
\begin{equation}
\lambda_J \sim \frac{c}{H} \cdot \frac{1}{\sqrt{\Omega_{\mathrm{mech}}|1+\wmech|}} \sim R_H \cdot (\text{geometric factor}).
\end{equation}
For $\Omega_{\mathrm{mech}} \ll 1$, the Jeans scale exceeds the Hubble radius, confirming that $\umech$ is smooth on all observable scales.
\end{proposition}

\subsection{Coupling to Photon-Baryon Fluid and Gravitational Potentials}
\label{sec:coupling}

The transient component $\umech$ couples to other cosmological fluids through the gravitational potential.

\subsubsection{Modified Poisson Equation}

In Fourier space, the perturbed Einstein equations give:
\begin{equation}
\label{eq:poisson_perturbed}
-k^2\Phi = 4\pi G a^2 \sum_i \bar{\rho}_i \delta_i,
\end{equation}
where the sum runs over all species (matter, radiation, $\umech$, etc.).

With $\umech$ included:
\begin{equation}
\boxed{-k^2\Phi = 4\pi G a^2 \left[ \rho_m \delta_m + \rho_r \delta_r + \umech \delta_{\mathrm{mech}} + \ldots \right]}
\end{equation}

Since $\delta_{\mathrm{mech}} \approx 0$ on sub-horizon scales (smooth component), $\umech$ contributes primarily through the background:
\begin{equation}
H^2 = \frac{8\pi G}{3}(\bar{\rho}_m + \bar{\rho}_r + \bar{\umech}/c^2 + \ldots).
\end{equation}

\subsubsection{Effects on Photon-Baryon Oscillations}

The photon-baryon fluid satisfies:
\begin{equation}
\ddot{\delta}_\gamma + \frac{\dot{R}}{1+R}\dot{\delta}_\gamma + c_s^2 k^2 \delta_\gamma = -k^2\Psi - \frac{4}{3}k^2\Phi + \frac{4\ddot{\Phi}}{3} + \ldots
\end{equation}

The $\umech$ component affects this equation through:
\begin{enumerate}
\item \textbf{Background effect:} Modified $H(a)$ changes the damping and oscillation rates.
\item \textbf{Metric source:} $\umech$ contributes to $\Phi$ and $\Psi$ (though negligibly if smooth).
\item \textbf{ISW effect:} Time-varying $\Phi$ generates late-time anisotropies.
\end{enumerate}

\subsubsection{Integrated Sachs-Wolfe Effect from $\umech$ Decay}

The ISW contribution to CMB temperature anisotropies is:
\begin{equation}
\label{eq:ISW}
\left(\frac{\Delta T}{T}\right)_{\mathrm{ISW}} = -2\int_0^{\tau_0} \dot{\Phi} \, d\tau.
\end{equation}

When $\umech$ decays (for $a > \ac$), it contributes to the time variation of $\Phi$. For a smooth component:
\begin{equation}
\dot{\Phi} \sim -\mathcal{H}\Phi - \frac{4\pi G a^2}{k^2}\dot{\bar{\umech}}\delta_{\mathrm{mech}} + \ldots
\end{equation}

Since $\delta_{\mathrm{mech}} \approx 0$ on relevant scales, the primary effect is through the background evolution modifying $\Phi(a)$ via the changed expansion history.

\begin{proposition}[ISW Contribution from $\umech$]
\label{prop:ISW_mech}
The decay of $\umech$ after $a = \ac$ produces an \textbf{early} ISW effect (if $\ac$ is during radiation or early matter era). The ISW contribution scales as:
\begin{equation}
\left(\frac{\Delta T}{T}\right)_{\mathrm{ISW}}^{\mathrm{mech}} \sim \int_{a_{\mathrm{mech,on}}}^{a_{\mathrm{mech,off}}} \dot{\Phi} \, d\tau \sim \Phi \cdot \Delta\left(\frac{\dot{a}}{a}\right),
\end{equation}
where the integral is over the epoch when $\umech$ is dynamically significant.

For $\fmech \sim 0.1$, this produces ISW contributions at the percent level on large scales ($\ell \lesssim 30$).
\end{proposition}


%==============================================================================
% PART 3: CMB POWER SPECTRUM MODIFICATIONS
%==============================================================================

\section{CMB Power Spectrum Modifications}
\label{sec:cmb_spectrum}

We now compute how $\umech$ modifies the CMB temperature and polarization power spectra.

\subsection{CMB Angular Power Spectrum Formalism}
\label{sec:cmb_formalism}

The CMB temperature angular power spectrum is:
\begin{equation}
\label{eq:Cl_TT}
C_\ell^{TT} = \frac{4\pi}{(2\ell+1)^2} \int_0^\infty \frac{dk}{k} \, \mathcal{P}_\mathcal{R}(k) \, |\Delta_T(k,\ell)|^2,
\end{equation}
where $\mathcal{P}_\mathcal{R}(k)$ is the primordial curvature power spectrum and $\Delta_T(k,\ell)$ is the temperature transfer function.

The transfer function encodes all physics between the primordial epoch and today:
\begin{equation}
\Delta_T(k,\ell) = \int_0^{\tau_0} d\tau \, S_T(k,\tau) \, j_\ell[k(\tau_0 - \tau)],
\end{equation}
where $S_T$ is the source function and $j_\ell$ are spherical Bessel functions.

\subsection{Effects of $\umech$ on the Transfer Function}
\label{sec:transfer_function}

\subsubsection{Background Effects}

The modified expansion history $H(a)$ affects the transfer function through:

\paragraph{1. Acoustic Peak Positions.}
The $n$-th acoustic peak is located at:
\begin{equation}
\label{eq:peak_position}
\ell_n \approx n\pi \frac{\dA(\zdec)}{\rs(\zdec)}.
\end{equation}

If $\umech$ reduces $\rs$ by $\Delta\rs$ while leaving $\dA$ approximately unchanged (valid for early-time modifications):
\begin{equation}
\frac{\Delta\ell_n}{\ell_n} \approx -\frac{\Delta\rs}{\rs} > 0.
\end{equation}

\begin{proposition}[Peak Shift from $\umech$]
\label{prop:peak_shift}
A positive $\umech$ (reducing $\rs$) shifts acoustic peaks to \textbf{higher} $\ell$ (smaller angular scales):
\begin{equation}
\boxed{\Delta\ell_n \approx \ell_n \cdot \frac{|\Delta\rs|}{\rs} > 0}
\end{equation}
For $|\Delta\rs/\rs| \sim 0.5\%$, the first peak ($\ell_1 \approx 220$) shifts by $\Delta\ell_1 \sim 1$.
\end{proposition}

\paragraph{2. Photon Diffusion Damping.}
The Silk damping scale is:
\begin{equation}
r_D^2 = \int_0^{\adec} \frac{da}{a^3 H n_e \sigma_T} \cdot \frac{R^2 + 16(1+R)/15}{6(1+R)^2},
\end{equation}
where $n_e$ is the electron density and $\sigma_T$ is the Thomson cross-section.

Modified $H(a)$ changes $r_D$:
\begin{equation}
\frac{\Delta r_D}{r_D} \approx -\frac{1}{2}\frac{\Delta H}{H} \approx -\frac{1}{4}\frac{\Omega_{\mathrm{mech}}}{E_{\mathrm{std}}^2} < 0.
\end{equation}

A larger $H$ (from positive $\umech$) reduces the damping scale, meaning more small-scale power survives. The damping tail ($\ell > 1000$) is enhanced.

\paragraph{3. Peak Height Ratios.}
The relative heights of acoustic peaks depend on the baryon density (through $R$) and the radiation driving force. $\umech$ modifies the driving through the metric perturbation:
\begin{equation}
\text{Driving} \propto \Phi + \Psi/3.
\end{equation}

During radiation domination with $\umech$, the potential decay is modified, changing the driving of acoustic oscillations. This primarily affects odd-even peak height ratios.

\subsubsection{Perturbation Effects}

Since $\umech$ is smooth ($c_s^2 = c^2$), direct perturbation effects are subdominant. The primary perturbation effect is through the modified potential evolution sourcing the ISW effect.

\subsection{Quantitative Estimates of CMB Modifications}
\label{sec:cmb_quantitative}

\begin{table}[H]
\centering
\caption{CMB Modifications from $\umech$ ($\fmech = 0.1$, $\ac = 10^{-3}$, $\sigmamech = 0.2$)}
\label{tab:cmb_modifications}
\begin{tabular}{@{}lcc@{}}
\toprule
\textbf{Observable} & \textbf{$\Lambda$CDM Value} & \textbf{Shift from $\umech$} \\
\midrule
Sound horizon $\rs$ & 144.4 Mpc & $-0.5$ Mpc ($-0.35\%$) \\
First peak $\ell_1$ & 220.0 & $+0.8$ ($+0.35\%$) \\
Peak spacing $\Delta\ell$ & 302 & $+1.1$ ($+0.35\%$) \\
Damping scale $\ell_D$ & 1500 & $+3$ ($+0.2\%$) \\
ISW amplitude (low $\ell$) & $\sim 1000 \mu\text{K}^2$ & $+30 \mu\text{K}^2$ ($+3\%$) \\
\bottomrule
\end{tabular}
\end{table}

\begin{proposition}[Detectability of $\umech$ Effects]
\label{prop:detectability}
For Paper 1's fiducial parameters, the CMB modifications are at the $0.3$--$0.5\%$ level for peak positions and $\sim 3\%$ for ISW amplitude. These are:
\begin{itemize}
\item \textbf{Below Planck sensitivity} for peak positions ($\sigma_{\ell_1} \sim 0.3$, i.e., $\sim 0.1\%$).
\item \textbf{Marginally detectable} for ISW amplitude enhancement.
\item \textbf{Degenerate} with other parameters (e.g., $\Omega_b h^2$, $n_s$).
\end{itemize}

Detection requires either:
\begin{enumerate}
\item Larger $\fmech$ parameters (tension with other constraints), or
\item Combined analysis with BAO, SNe, and other probes.
\end{enumerate}
\end{proposition}

\subsection{Summary: CMB Constraints on $\umech$}
\label{sec:cmb_summary}

The CMB constrains the $\umech$ parameter space through:
\begin{enumerate}
\item \textbf{Acoustic peak positions:} Constrain integrated $\Omega_{\mathrm{mech}}$ effect on $\rs$.
\item \textbf{Damping tail:} Constrain $\umech$ amplitude during recombination epoch.
\item \textbf{ISW effect:} Constrain decay rate of $\umech$ after peak.
\item \textbf{BBN consistency:} Constrain $\umech$ during nucleosynthesis ($a \sim 10^{-9}$).
\end{enumerate}

Current Planck data are consistent with no $\umech$ detection, placing upper limits of roughly:
\begin{equation}
\fmech \lesssim 0.1 \quad \text{for } \ac \lesssim \adec.
\end{equation}


%==============================================================================
% PART 4: MATTER POWER SPECTRUM AND S_8 TENSION
%==============================================================================

\section{Matter Power Spectrum and the $S_8$ Tension}
\label{sec:matter_spectrum}

\subsection{Matter Power Spectrum with $\umech$}
\label{sec:Pk_mech}

The matter power spectrum today is:
\begin{equation}
P(k) = A_s k^{n_s} T^2(k) D^2(a=1),
\end{equation}
where $A_s$ is the primordial amplitude, $n_s$ is the spectral index, $T(k)$ is the transfer function, and $D(a)$ is the growth function.

The $\umech$ component modifies $P(k)$ through:
\begin{enumerate}
\item \textbf{Jeans cutoff:} The BEC healing length $\xi \sim 1$ kpc introduces a small-scale cutoff (addressed in Paper 3).
\item \textbf{Growth function:} $\umech$ modifies the growth rate during its active epoch.
\end{enumerate}

\subsection{Growth Function with $\umech$}
\label{sec:growth_function}

The linear growth function $D(a)$ satisfies:
\begin{equation}
\label{eq:growth_ODE}
D'' + \left(2 + \frac{H'}{H}\right)D' - \frac{3}{2}\Omega_m(a) D = 0,
\end{equation}
where primes denote derivatives with respect to $\ln a$, and:
\begin{equation}
\Omega_m(a) = \frac{\Omm a^{-3}}{E^2(a)} = \frac{\Omm a^{-3}}{\Omr a^{-4} + \Omm a^{-3} + \OmL + \Omega_{\mathrm{mech}}(a)}.
\end{equation}

\subsubsection{Effect of $\umech$ on Growth}

During the $\umech$-dominated epoch (if any), $\Omega_m(a)$ is suppressed:
\begin{equation}
\Omega_m(a) \approx \frac{\Omm a^{-3}}{\Omega_{\mathrm{mech}}(a)} < \Omega_m^{\mathrm{std}}(a).
\end{equation}

This \textbf{suppresses} the gravitational source term $\frac{3}{2}\Omega_m D$, reducing the growth rate:
\begin{equation}
\frac{D_{\mathrm{mech}}}{D_{\mathrm{std}}} < 1.
\end{equation}

\begin{proposition}[Growth Suppression from $\umech$]
\label{prop:growth_suppression}
A positive $\umech$ component suppresses matter growth during its active epoch. The suppression is:
\begin{equation}
\boxed{\frac{\Delta D}{D} \approx -\frac{3}{4}\int_{\ln a_i}^{\ln a_f} \frac{\Omega_{\mathrm{mech}}(a)}{E^2(a)} \, d\ln a < 0}
\end{equation}
For $\fmech = 0.1$, $\sigmamech = 0.2$, $\ac = 10^{-3}$, the growth suppression is of order $0.1$--$1\%$.
\end{proposition}

\subsection{Effect on $\sigma_8$ and $S_8$}
\label{sec:sigma8}

The amplitude of matter fluctuations at scale $8 h^{-1}$ Mpc is:
\begin{equation}
\sigma_8^2 = \frac{1}{2\pi^2}\int_0^\infty k^2 P(k) W^2(kR_8) \, dk,
\end{equation}
where $W$ is a top-hat window function and $R_8 = 8 h^{-1}$ Mpc.

Since $\umech$ suppresses growth:
\begin{equation}
\sigma_8^{\mathrm{mech}} < \sigma_8^{\mathrm{std}}.
\end{equation}

The derived parameter $S_8 = \sigma_8\sqrt{\Omm/0.3}$ is similarly reduced:
\begin{equation}
S_8^{\mathrm{mech}} < S_8^{\mathrm{std}}.
\end{equation}

\subsection{The $S_8$ Tension}
\label{sec:S8_tension}

Current measurements show a tension:
\begin{align}
S_8^{\mathrm{Planck}} &= 0.834 \pm 0.016, \\
S_8^{\mathrm{Weak\,Lensing}} &= 0.759 \pm 0.024.
\end{align}

The weak lensing value is $\sim 3\sigma$ \textbf{lower} than Planck.

\begin{proposition}[BEC Framework and the $S_8$ Tension]
\label{prop:S8_tension}
The $\umech$ component suppresses growth, reducing $\sigma_8$ and $S_8$. This is the \textbf{correct direction} to alleviate the $S_8$ tension:
\begin{equation}
S_8^{\mathrm{mech}} < S_8^{\mathrm{Planck}} \quad \text{(moves toward weak lensing value)}.
\end{equation}

For $\fmech \sim 0.1$, the suppression is $|\Delta S_8/S_8| \sim 0.5$--$1\%$, which is \textbf{insufficient} to fully resolve the $\sim 10\%$ tension.

However, if $\umech$ has larger amplitude or persists longer (broader $\sigmamech$), the suppression could be significant. The framework provides a physical mechanism (early dark energy-like behavior) that qualitatively addresses both the $H_0$ and $S_8$ tensions in the correct directions.
\end{proposition}

\subsection{Combined Constraints: $H_0$ vs. $S_8$}
\label{sec:combined_constraints}

There is a degeneracy between $H_0$ and $S_8$ modifications:
\begin{itemize}
\item \textbf{Increasing $H_0$} (via reduced $\rs$) typically \textbf{increases} $S_8$ (more time for growth at fixed physical scales).
\item \textbf{Decreasing $S_8$} (via growth suppression) typically \textbf{decreases} inferred $H_0$.
\end{itemize}

The $\umech$ component acts in a specific way:
\begin{enumerate}
\item \textbf{Before recombination:} Increases $H$, reduces $\rs$, increases inferred $H_0$.
\item \textbf{During matter era:} Suppresses $\Omega_m(a)$, reduces growth, decreases $S_8$.
\end{enumerate}

\begin{theorem}[Parameter Space for Simultaneous Tension Resolution]
\label{thm:parameter_space}
To simultaneously resolve both tensions requires:
\begin{enumerate}
\item $\ac \sim \adec$ or slightly earlier (maximize impact on $\rs$).
\item $\sigmamech \gtrsim 0.5$ (extend impact into matter era for growth suppression).
\item $\fmech \sim 0.3$--$0.5$ (sufficient amplitude).
\end{enumerate}

However, such parameters are strongly constrained by:
\begin{itemize}
\item BBN ($\umech$ during nucleosynthesis changes light element abundances).
\item CMB damping tail (high-$\ell$ Planck data).
\item BAO measurements (constrain $\rs$ independently).
\end{itemize}

A full MCMC analysis with Boltzmann code integration is required to determine the viable parameter space.
\end{theorem}


%==============================================================================
% SUMMARY AND CONCLUSIONS
%==============================================================================

\section{Summary of Mathematical Results}
\label{sec:summary}

\begin{proposition}[Summary: Paper 4 Key Results]
\label{prop:summary_paper4}
We have derived the following mathematical foundations for CMB and matter spectrum analysis in the BEC framework:

\textbf{Part 1: Sound Horizon}
\begin{enumerate}
\item The sound horizon is reduced by $\umech$:
\begin{equation}
\Delta\rs = -\int_0^{\adec} \frac{c_s(a)}{a^2 H_0 E_{\mathrm{std}}} \cdot \frac{\Omega_{\mathrm{mech}}}{2E_{\mathrm{std}}^2} \, da < 0.
\end{equation}

\item For fiducial parameters ($\fmech = 0.1$, $\ac = 10^{-3}$, $\sigmamech = 0.2$):
\begin{equation}
\frac{|\Delta\rs|}{\rs} \approx 0.3\text{--}0.5\% \quad \text{(insufficient for $H_0$ tension)}.
\end{equation}

\item Resolving $H_0$ tension requires $|\Delta\rs/\rs| \approx 7$--$8\%$, demanding more aggressive parameters.
\end{enumerate}

\textbf{Part 2: Perturbations}
\begin{enumerate}
\item Perturbation evolution equations in conformal Newtonian gauge:
\begin{align}
\dot{\delta}_{\mathrm{mech}} &= -(1+\wmech)(\theta_{\mathrm{mech}} - 3\dot{\Phi}) - 3\mathcal{H}(c_s^2 - \wmech)\delta_{\mathrm{mech}}, \\
\dot{\theta}_{\mathrm{mech}} &= -\mathcal{H}(1-3\wmech)\theta_{\mathrm{mech}} + \frac{c_s^2 k^2}{1+\wmech}\delta_{\mathrm{mech}} + k^2\Psi.
\end{align}

\item Sound speed $c_s^2 = c^2$ (relativistic); $\umech$ is a smooth component with infinite Jeans length at the peak.

\item Coupling to other fluids is primarily through background $H(a)$ modification.
\end{enumerate}

\textbf{Part 3: CMB Spectrum}
\begin{enumerate}
\item Peak positions shift to higher $\ell$ by $\sim |\Delta\rs/\rs|$.
\item Damping tail enhanced by faster expansion (smaller $r_D$).
\item ISW contribution from $\umech$ decay at few-percent level on large scales.
\item Current Planck constraints: $\fmech \lesssim 0.1$ for $\ac \lesssim \adec$.
\end{enumerate}

\textbf{Part 4: Matter Spectrum and $S_8$}
\begin{enumerate}
\item Growth function suppressed during $\umech$-active epoch.
\item $\sigma_8$ and $S_8$ are reduced---correct direction for $S_8$ tension.
\item Fiducial parameters give $|\Delta S_8/S_8| \sim 0.5$--$1\%$ (insufficient for full resolution).
\item Framework provides mechanism to address both $H_0$ and $S_8$ tensions qualitatively.
\end{enumerate}

\end{proposition}


%==============================================================================
% APPENDIX: DETAILED INTEGRALS
%==============================================================================

\appendix

\section{Detailed Integral Evaluations}
\label{app:integrals}

\subsection{Sound Horizon Integral with Gaussian $\Omega_{\mathrm{mech}}$}

Starting from:
\begin{equation}
\Delta\rs = \frac{c}{2H_0\sqrt{3}} \int_0^{\adec} \frac{a^4}{\sqrt{\Omr}^3 \sqrt{1+R_0 a}} \cdot \fmech e^{-\ln^2(a/\ac)/(2\sigmamech^2)} \, da.
\end{equation}

Substituting $x = \ln(a/\ac)$, $a = \ac e^x$, $da = \ac e^x dx$:
\begin{equation}
\Delta\rs = \frac{c \fmech \ac^5}{2H_0\sqrt{3\Omr^3}} \int_{-\infty}^{\ln(\adec/\ac)} \frac{e^{5x}}{\sqrt{1+R_0\ac e^x}} \cdot e^{-x^2/(2\sigmamech^2)} \, dx.
\end{equation}

For narrow Gaussians ($\sigmamech \ll 1$), the integral is dominated near $x = 0$:
\begin{equation}
\Delta\rs \approx \frac{c \fmech \ac^5}{2H_0\sqrt{3\Omr^3(1+R_0\ac)}} \cdot \sigmamech\sqrt{2\pi}.
\end{equation}

\subsection{Growth Function Integral}

The growth suppression integral:
\begin{equation}
\frac{\Delta D}{D} \approx -\frac{3}{4}\int_{\ln a_i}^{\ln a_f} \frac{\fmech e^{-\ln^2(a/\ac)/(2\sigmamech^2)}}{E^2(a)} \, d\ln a.
\end{equation}

During matter domination, $E^2(a) \approx \Omm a^{-3}$:
\begin{equation}
\frac{\Delta D}{D} \approx -\frac{3\fmech}{4\Omm} \int a^3 e^{-\ln^2(a/\ac)/(2\sigmamech^2)} \, d\ln a.
\end{equation}

With $x = \ln(a/\ac)$:
\begin{equation}
\frac{\Delta D}{D} \approx -\frac{3\fmech \ac^3}{4\Omm} \int e^{3x} e^{-x^2/(2\sigmamech^2)} \, dx = -\frac{3\fmech \ac^3}{4\Omm} \cdot \sigmamech\sqrt{2\pi} \cdot e^{9\sigmamech^2/2}.
\end{equation}

For $\fmech = 0.1$, $\ac = 10^{-3}$, $\sigmamech = 0.2$, $\Omm = 0.31$:
\begin{equation}
\frac{\Delta D}{D} \approx -\frac{3 \times 0.1 \times 10^{-9}}{4 \times 0.31} \times 0.2\sqrt{2\pi} \times e^{0.18} \approx -10^{-10}.
\end{equation}

This is negligible because $\ac^3 \ll 1$. To get significant growth suppression, either $\ac$ must be larger (later peak) or $\sigmamech$ must be much larger (broader peak extending into matter era).


%==============================================================================
% END OF MATHEMATICAL DERIVATIONS
%==============================================================================
